% ============================================================================
% Trabajo de grado — Universidad Distrital Francisco Jose de Caldas
% Ingenieria de Sistemas
%
% Compilar:  cd paper && latexmk main.tex
% ============================================================================
\documentclass[12pt,letterpaper,oneside]{report}

% ====== Preambulo compartido — tesis TCLab (LQR + RL) ======
% Estandar definido en .claude/rules/thesis-format.md

% --- Idioma (primero: condiciona el resto) ---
\usepackage[spanish,es-tabla,es-noquoting]{babel}
\usepackage{csquotes}

% --- Diseno de pagina ---
\usepackage[left=3cm,right=2.5cm,top=2.5cm,bottom=2.5cm]{geometry}
\usepackage{setspace}
\onehalfspacing
\usepackage{fancyhdr}
\pagestyle{fancy}
\fancyhf{}
\fancyfoot[C]{\thepage}
\renewcommand{\headrulewidth}{0pt}

% --- Tipografia (XeLaTeX) ---
\usepackage{fontspec}
\setmainfont{Latin Modern Roman}
\usepackage{microtype}

% --- Titulos ---
\usepackage{titlesec}
\usepackage[title]{appendix}

% --- Matematicas ---
\usepackage{amssymb, amsmath, amsfonts, mathtools}
\usepackage{amsthm}
\theoremstyle{plain}
\newtheorem{teorema}{Teorema}[chapter]
\newtheorem{proposicion}[teorema]{Proposición}
\newtheorem{lema}[teorema]{Lema}
\newtheorem{corolario}[teorema]{Corolario}
\theoremstyle{definition}
\newtheorem{definicion}[teorema]{Definición}
\newtheorem{supuesto}[teorema]{Supuesto}
\DeclareMathOperator*{\argmax}{arg\,max}
\DeclareMathOperator*{\argmin}{arg\,min}
\DeclareMathOperator{\tr}{tr}
\newcommand{\norm}[1]{\left\lVert #1 \right\rVert}
\newcommand{\transpose}{^{\mathsf{T}}}

% --- Tablas ---
\usepackage{array, booktabs, makecell}
\usepackage{siunitx}
\sisetup{output-decimal-marker={,}}
\usepackage[flushleft]{threeparttable}
\usepackage{tabularx, rotating}
\usepackage{tabularray}
\UseTblrLibrary{booktabs, siunitx}

% --- Figuras ---
\usepackage{graphicx, subcaption}
\usepackage{float}
\usepackage{caption}
\captionsetup{font=small, labelfont=bf, justification=justified}

% --- Diagramas de bloques ---
\usepackage{tikz}
\usetikzlibrary{shapes, arrows.meta, positioning, calc}

% --- Codigo fuente ---
\usepackage{listings}
\lstset{basicstyle=\ttfamily\small, breaklines=true, frame=single,
        showstringspaces=false, captionpos=b}

% --- Listas ---
\usepackage{enumitem}

% --- Bibliografia: biblatex + biber, estilo IEEE ---
\usepackage{xurl}
\usepackage{xcolor}
\definecolor{citationcolor}{RGB}{0, 90, 160}
\usepackage[backend=biber, style=ieee, sorting=none, maxbibnames=99]{biblatex}
\addbibresource{Bibliography_base.bib}

% --- Notas al pie ---
\interfootnotelinepenalty=10000

% --- hyperref (penultimo) ---
\usepackage[breaklinks, colorlinks=true,
            linkcolor=citationcolor, citecolor=citationcolor,
            urlcolor=citationcolor]{hyperref}

% --- cleveref (despues de hyperref) ---
\usepackage[spanish,nameinlink]{cleveref}


% ---- Datos del trabajo (actualizar) ----
\newcommand{\titulotesis}{Control de temperatura en la planta TCLab mediante LQR
combinado con aprendizaje por refuerzo}
\newcommand{\autortesis}{[Nombre del autor]}
\newcommand{\codigoestudiante}{[Codigo estudiantil]}
\newcommand{\directortesis}{[Nombre del director]}
\newcommand{\aniotesis}{2026}

\begin{document}

% ============================ PORTADA ============================
\begin{titlepage}
\thispagestyle{empty}
\centering

{\large\bfseries \titulotesis \par}
\vspace{2.5cm}

{\normalsize \autortesis \par}
{\normalsize Código: \codigoestudiante \par}
\vspace{2.5cm}

{\normalsize Trabajo de grado para optar al título de\\
Ingeniero de Sistemas \par}
\vspace{2cm}

{\normalsize Director:\\ \directortesis \par}
\vfill

{\normalsize
Universidad Distrital Francisco José de Caldas\\
Facultad de Ingeniería\\
Ingeniería de Sistemas\\
Bogotá D.C.\\
\aniotesis \par}

\end{titlepage}

% ==================== PAGINAS PRELIMINARES ====================
\pagenumbering{roman}
\setcounter{page}{2}

\chapter*{Resumen}
\addcontentsline{toc}{chapter}{Resumen}

[Pendiente. Se redacta al final, cuando existan resultados. Debe respetar el límite del
reglamento de trabajos de grado del programa (INV-5).]

\vspace{1em}
\noindent\textbf{Palabras clave:} control óptimo, LQR, aprendizaje por refuerzo, TCLab,
identificación de sistemas.

\chapter*{Abstract}
\addcontentsline{toc}{chapter}{Abstract}

[Pending.]

\vspace{1em}
\noindent\textbf{Keywords:} optimal control, LQR, reinforcement learning, TCLab,
system identification.


\tableofcontents
\listoffigures
\listoftables

% ======================== CUERPO ========================
\cleardoublepage
\pagenumbering{arabic}
\setcounter{page}{1}

\chapter{Introducción}
\label{cap:introduccion}

\section{Planteamiento del problema}
[Pendiente]

\section{Justificación}
[Pendiente]

\section{Objetivos}

\subsection{Objetivo general}
[Pendiente — se define tras \texttt{/discover interview}]

\subsection{Objetivos específicos}
% Los objetivos especificos son un contrato con el jurado: el capitulo de
% conclusiones debe cerrar cada uno explicitamente.
\begin{enumerate}
    \item [Pendiente]
\end{enumerate}

\section{Alcance y limitaciones}
[Pendiente]

\section{Estructura del documento}
[Pendiente]

\chapter{Marco teórico}
\label{cap:marco_teorico}

% Solo lo que el trabajo usa. Un marco teorico enciclopedico y desconectado
% de las decisiones tomadas es un hallazgo del evaluador de dominio.

\section{Representación en espacio de estados}
[Pendiente]

\section{Identificación de sistemas}
[Pendiente]

\section{Regulador lineal cuadrático}
% Fijar aqui la convencion de notacion y sostenerla en todo el documento (INV-7).
% Atencion a la colision: Q es la matriz de peso del LQR, y Q1/Q2 son las potencias
% de los calentadores en la nomenclatura del TCLab. Elegir una convencion y declararla.
[Pendiente]

\section{Aprendizaje por refuerzo}
[Pendiente]

\section{Notación}
[Tabla de símbolos: qué significa cada uno y en qué unidades.]

\chapter{Estado del arte}
\label{cap:estado_del_arte}

\section{Control óptimo aplicado a procesos térmicos}
[Pendiente]

\section{Aprendizaje por refuerzo en control de procesos}
[Pendiente]

\section{Combinación de control clásico y aprendizaje}
[Pendiente]

\section{Posicionamiento de este trabajo}
[Pendiente — en qué se diferencia de lo revisado]

\chapter{Metodología}
\label{cap:metodologia}

\section{Plataforma experimental}
% Descripcion del TCLab: estructura MIMO 2x2, actuadores, sensores, rangos.
[Pendiente]

\section{Protocolo de identificación}
% Senal de excitacion, amplitud, duracion, punto de operacion, tiempo de muestreo.
% Corridas de identificacion y de validacion SEPARADAS (INV-25).
[Pendiente]

\section{Diseño del controlador}
[Pendiente]

\section{Formulación del componente de aprendizaje}
[Pendiente]

\section{Protocolo experimental}
% Perfil de setpoints identico para todos los controladores comparados (INV-24).
% Repeticiones, semillas, tiempo de enfriamiento, registro de temperatura ambiente.
% Criterio de exito declarado ANTES de correr los experimentos.
[Pendiente]

\section{Métricas de evaluación}
% Con unidades explicitas (INV-4).
[Pendiente]

\chapter{Implementación}
\label{cap:implementacion}

\section{Arquitectura de software}
[Pendiente]

\section{Capa de adquisición y seguridad}
% Apagado garantizado de calentadores, saturacion, corte por sobretemperatura (INV-20, INV-21).
[Pendiente]

\section{Lazo de control en tiempo real}
% Temporizacion por deadline absoluto, registro de instantes reales, desbordes de periodo.
[Pendiente]

\section{Entorno de simulación y entrenamiento}
[Pendiente]

\chapter{Resultados}
\label{cap:resultados}

\section{Identificación del modelo}
% Ajuste sobre datos de validacion independientes, con el criterio y su valor.
[Pendiente]

\section{Desempeño en simulación}
[Pendiente]

\section{Desempeño en hardware}
% Media +/- desviacion estandar sobre N corridas, con N declarado (INV-4).
[Pendiente]

\section{Comparación entre controladores}
% Condiciones identicas (INV-24). Ninguna diferencia se declara relevante si es
% menor que la dispersion entre corridas de la misma configuracion.
[Pendiente]

\chapter{Discusión}
\label{cap:discusion}

\section{Interpretación de los resultados}
[Pendiente]

\section{Brecha entre simulación y hardware}
[Pendiente]

\section{Limitaciones}
% Reconocer las limitaciones suma ante el jurado; ocultarlas resta.
[Pendiente]

\chapter{Conclusiones y trabajo futuro}
\label{cap:conclusiones}

\section{Cumplimiento de los objetivos}
% Cerrar EXPLICITAMENTE cada objetivo especifico declarado en la introduccion.
% El jurado los verifica uno a uno.
[Pendiente]

\section{Conclusiones}
[Pendiente]

\section{Trabajo futuro}
[Pendiente]


% ===================== REFERENCIAS =====================
\cleardoublepage
\printbibliography[heading=bibintoc, title={Referencias}]

% ======================= ANEXOS =======================
\begin{appendices}
\chapter{Anexos}
\label{cap:anexos}

\section{Configuración de los experimentos}
[Pendiente — contenido de \texttt{config/}]

\section{Código fuente relevante}
[Pendiente]

\end{appendices}

\end{document}
